%%
%% Author: shoupu_wan
%% 8/22/18
%%

\chapter{Json prettifier}\label{ch:json-prettifier}

\section{The problem statement}

The requirement is simply as ``Given a string as a Json blob, make it pretty, well-formatted text''.

\section{Outline of the algorithm}

The algorithm consists of two steps.
Step one, parse Json data;
Step two, reformat the Json data by adding newlines, indent spaces, \textit{etc.}

\section{Implementation}

Divide and conquer.
The input string may be a blob of strings with no newlines or spaces in between.
Or the input may be `almost' well-formated Json files, off by just a few spaces or
newlines---this requires minimal `style-compliance' formatting.
Or the input may be invalid Json, with some missing part or extra part.

This inspires us to validate and parse the input string into an in-memory data structure
first and then print the structure as pretty string.

\autoref{lst:JsonPrettifier.py},  \autoref{lst:JsonFormatter.py}, and \autoref{lst:JsonParser.py}.

\begin{minipage}{\linewidth}
    \lstinputlisting[label={lst:JsonPrettifier.py},firstline=10,lastline=31,
    caption={Json prettifier code that invokes the modular constructs of \code{JsonParser} and
    \code{JosnFormatter} to prettify input Json blob data.}]
    {case-study/my-json/JsonPrettifier.py}
\end{minipage}

\begin{minipage}{\linewidth}
    \lstinputlisting[label={lst:JsonFormatter.py},firstline=10,lastline=31,
    caption={Json formatter.
    It take an in-memory object reprensetation of Json data, print it into well-formatted string.}]
    {case-study/my-json/JsonFormatter.py}
\end{minipage}

\begin{minipage}{\linewidth}
    \lstinputlisting[label={lst:JsonParser.py},firstline=10,lastline=31,
    caption={Json formatter.
    It take an in-memory object reprensetation of Json data, print it into well-formatted string.}]
    {case-study/my-json/JsonParser.py}
\end{minipage}
